He felt them before they arrived — not the lords and captains filing into the Great Hall for the war council, whose footsteps any ordinary man could have heard well enough through stone, but the others, the ones still three days’ hard march north, moving down the coast in numbers that no living scout had yet managed to count properly, because no living scout who got close enough to count them ever came back to report the number.

Bran had stopped, some weeks past, trying to explain to people what it felt like to sense an army the way another man felt weather coming — a pressure behind the eyes, a wrongness in the way the crows overhead had begun flying in patterns that meant nothing to anyone who hadn’t spent months learning to read a raven’s flight the way other men read a raven’s message. He had given up trying to make Sansa understand it properly, or Jon, or even Sam, who came closest of anyone still breathing to grasping what he’d become, and had settled instead for simply telling people numbers, distances, timing, and letting them assume the knowing came easier to him than it did.

It did not come easier. It came at a cost he had not yet found the words to explain to anyone, least of all himself.

\sceneBreak

“Forty thousand,” Bran said, once Podrick had wheeled him to his place at the map table and the hall’s murmur had quieted enough for his voice to carry. “Perhaps a few thousand more or less — the number shifts, because the number includes everything they’ve raised along the way, every village between Eastwatch and here that didn’t get its dead burned properly before the host reached it. They move faster than any living army could, because they don’t tire and they don’t stop for anything except open water deep enough to trouble even them, and the coast route they’re taking has almost none of that between here and the Wall’s ruin.”

“Three days,” Jon said, and it was not a question. He had asked Bran the same thing every day for a week now, and Bran understood, watching his brother’s face in the candlelight, that the asking was less about the answer than about needing to hear it said aloud by someone, the way a man touches a wound repeatedly not because he doubts it’s there but because some part of him needs to keep confirming it hasn’t somehow healed itself when he wasn’t looking.

“Three days. Perhaps two, if the weather favors them, which it generally seems to.” Bran let his eyes move over the map spread across the table — old vellum, Winterfell’s walls and the wolfswood and the kingsroad marked in a careful hand generations dead now, utterly inadequate to the thing it was being asked to help defeat. “They don’t march in ranks the way your soldiers do, Jon. I want that understood plainly, because I think some of the captains in this hall are still picturing something closer to an ordinary army, and it isn’t that. It moves like one thing. Like a hand, closing.”

“You’re saying it’s commanded.” Grey Worm’s voice, flat and precise, from further down the table. “Not merely herded, or driven. Commanded, the way a general commands.”

“I’m saying I’ve watched it change direction all at once, every wight in a column three miles long, without any signal passing between them that I could see or hear, the way your own soldiers might respond to a horn or a raised banner.” Bran held Grey Worm’s gaze, and found, in the Unsullied commander’s stillness, someone who understood immediately what most of the hall would need considerably longer to accept. “Whatever the Night King is, he isn’t only a monster leading monsters. He’s a mind, holding thousands of bodies the way you hold your own two hands, and I don’t think any of us should walk into this battle believing we’re facing something that can be outsmarted merely because it used to be a man once, a very long time ago.”

The hall had gone very quiet, and Bran felt, beneath the quiet, the particular fear of soldiers who had spent their whole lives training to fight other men and were only now, this late, beginning to understand that the thing coming for them had never once, in eight thousand years, needed to learn anything a living general could counter.

\sceneBreak

“Then we don’t try to outsmart it,” Jaime said, into the silence, and Bran watched heads turn toward the Kingslayer with the particular wary attention Jaime had not yet managed to fully outlive in this hall. “You don’t win a siege against a mind that doesn’t tire by being clever, Bran. You win it — or you survive it, which I begin to suspect is the more honest goal here — by making the ground itself do work no cleverness can undo. Trenches, deep enough that nothing climbs out of them once they’re burning. Choke points that force the numbers to matter less than the width of the gap they’re funneled through. I’ve besieged cities. I’ve been besieged in them. The dead may not tire, but fire doesn’t care whether the thing it’s burning gets tired first.”

“The wildfire,” Tyrion said, “assuming Sansa’s stores and what Winterfell’s own maester can improvise between them prove equal to numbers this large, which I confess worries me rather more than anything Bran’s just told us about coordinated wights.”

“They’ll do,” Sansa said, with the particular flatness of a woman who had already run the numbers three times before this meeting and did not intend to relitigate them for the hall’s comfort. “Enough for the outer trenches, and enough held back for the second line, if the first one fails to hold. I’ll not pretend it’s enough for a siege that runs longer than one night, but Bran tells us we haven’t the luxury of planning for one that runs longer than that regardless.”

“And the dragons?” This from Davos, practical as ever, eyes on Daenerys rather than on Bran.

“Will do more good in the air than tethered to any one position,” Daenerys said. “But not blindly. Not charging every shadow that moves, the way I understand some of your captains are hoping. Rhaegal and Drogon fly where Jon and I judge the weight of the fight actually needs breaking, and no further, because I’ve buried enough soldiers who died charging bravely at the wrong moment to know courage and strategy aren’t always the same virtue.” Her eyes found Bran’s, briefly, some old wariness in them that had not fully eased since her arrival, though it had softened considerably from where it started. “You said ‘a mind that doesn’t tire.’ Does that include whatever it’s riding? The dragon it turned?”

“I don’t know,” Bran said honestly, which cost him something to admit in a hall full of people who had come to expect certainty from him whether or not he’d earned it. “I’ve seen it in fragments, the way I see most things that haven’t happened yet or are happening too far away for the trees to show clearly. I believe it tires less than a living dragon would, and feels less, which may be worse for you rather than better, since a living dragon can be frightened off and this one, I think, cannot.”

\sceneBreak

Jorah Mormont spoke for the Dothraki, when the discussion turned to them, and Bran watched something old and complicated move behind the knight’s weathered face — a man arguing, in his own home country, for people who were not his own blood but had become, somewhere across two years and a hard winter’s march, something closer to a charge he meant to bring home alive if he could manage it.

“Not a charge,” Jorah said, when Ser Brynden’s old cavalry doctrine got raised by one of the Northern captains, some instinct toward a single decisive strike that Bran understood, watching it, came from men who had never fought anything that didn’t eventually break and run. “They don’t break. They don’t run. A charge into that only spends men for nothing. Let them ride the flanks instead — hit where the line’s thinnest, wheel away before it closes around them, hit again somewhere else. It costs discipline these riders aren’t accustomed to keeping, but it costs fewer lives than one glorious ride into a wall of dead that never once considers retreating.”

“They won’t like being told to run rather than charge,” Tormund said, from further down the table, arms crossed, the particular grim amusement of a man who had spent a lifetime among people with more pride than sense. “Ironborn, Dothraki, doesn’t much matter — proud folk hear ‘wheel away’ and hear ‘coward,’ half the time, whatever sense sits underneath the order.”

“Then we’d best make certain they understand the difference before the dead are close enough to make the lesson cost blood instead of pride,” Jon said, and something in the set of his jaw told Bran his brother had already spent considerable thought on exactly this problem, in whatever hours he wasn’t spending awake and afraid in the library or the godswood. “Theon.”

Theon looked up from where he’d been sitting quiet near the hall’s edge, some old habit of making himself smaller than the room required still visible in how he held his shoulders. “My lord?”

“The godswood approach. I want it held, not merely watched — Bran means to be there when this begins, and whatever else happens tonight, that ground doesn’t fall while he’s on it.” Jon’s eyes found Bran’s briefly, some old brotherly worry still visible beneath the King in the North’s careful command of the room. “Can you hold it?”

“I can hold it,” Theon said, and there was no performance in the saying of it, only a plain, settled certainty that Bran — who had, after all, more reason than anyone in this hall to weigh exactly what Theon Greyjoy’s word was worth — found he believed without reservation.

\sceneBreak

The council ran long past what any of them had planned for, numbers and positions and contingencies traded back and forth until even Tyrion’s wit had worn thin with exhaustion, and Bran sat through most of it saying less than the hall expected of him, because the part of this he understood most clearly was also the part he had the least language left to share with anyone still fully living.

He knew, in a way none of them yet did, that the godswood was not merely a defensible position. He knew that whatever happened there would cost him something the old Raven had warned him about long before he understood the warning was meant literally — that he would fly, and that flying would cost him the walking he’d already lost, and considerably more besides. He had not told Jon this. He had not told Sansa, or Sam, or anyone, because he had come to understand, these past weeks, that some knowledge only became a burden to the people who loved him once he’d actually spoken it aloud, and he found he could not yet bring himself to spend Jon’s remaining peace on a fear that would arrive soon enough regardless of whether Jon spent the intervening days dreading it.

When Podrick wheeled him back toward the godswood that night, past a castle settling into whatever fitful sleep it could manage before the numbers Bran had given them stopped being numbers and became instead the thing itself, arriving at last out of the dark, Bran found himself thinking, not for the first time, that there was a particular loneliness in being the only person in a hall full of frightened, brave people who already knew, with something closer to certainty than anyone else in that room would have permitted themselves, exactly how much this battle was going to cost him personally before it cost anyone else anything at all.